\subsection{Identidad y ciclo de vida de conductores externos}
\label{sec:conductores-externos}

Los conductores de transportistas externos ($\approx$\,160, de 10 empresas) no pertenecen a la dotación de Puelche.

\begin{itemize}
  \item Antes de entregar una ruta, la empresa transportista declara la identidad del conductor y su vínculo con ella.
  \item En la Etapa 1, despacho registra esa declaración en su función operacional.
  \item Desde la Etapa 2, el representante de la empresa puede confirmarla en el portal de transportistas.
  \item El responsable de despacho aprueba la asignación conductor--vehículo--turno; el alta inicial requiere conexión para verificar un OTP de un solo uso.
  \item El identificador del conductor y el del transportista acompañan cada \texttt{EntregaRegistrada}.
  \item La autorización se limita a la ruta asignada y a M6, M7 y M8 mediante rol y atributos de turno, sitio y dispositivo.
\end{itemize}

La aplicación conserva localmente una asignación firmada para el turno de hasta 14 horas y permite capturar eventos mientras no exista red.

\begin{itemize}
  \item La caché local de identidad, de solo lectura y TTL de 24 horas, no acorta por sí sola la vigencia de esa asignación de terreno; tampoco la renueva.
  \item Una persona reemplazante que no fue dada de alta no obtiene una nueva identidad sin conexión.
  \item En ese caso Operaciones reasigna el turno a un conductor ya habilitado y registra la excepción; la rotación no autoriza compartir el OTP ni la cuenta del conductor previo.
  \item Para que esa reasignación sea posible durante un corte de 24 horas, el acuerdo operacional con cada transportista le exige mantener enrolados, antes de cada turno, conductores suplentes suficientes para cubrir sus rutas del día siguiente.
  \item El manifiesto firmado que Keycloak renueva al menos cada hora los incluye.
  \item Si aun así no hay un suplente habilitado, la ruta no sale con una persona no identificada: Operaciones la reasigna a un camión con conductor habilitado o la reprograma.
  \item Este caso se ensaya en la prueba de corte de 24 horas antes del hito H5 del Formulario E-25.
\end{itemize}

Al finalizar el turno expiran la asignación y sus permisos.

\begin{itemize}
  \item Si se denuncia pérdida de dispositivo o baja anticipada, Keycloak revoca el acceso conectado y el equipo borra los datos al recuperar conexión.
  \item La revocación inmediata de RT-12.07 (Bases Técnicas Transversales, cap. 12, p. 25) se cumple en todo lo que depende de la solución.
  \item Al recibir la denuncia, Keycloak revoca la sesión y las credenciales, la API rechaza toda operación posterior del dispositivo y la sincronización rechaza los eventos firmados después de la hora de revocación, que quedan en cuarentena para revisión.
  \item Un equipo sin señal no puede borrarse antes de que se conecte.
  \item Para ese intervalo, la aplicación exige el PIN personal en cada sesión, cifra los datos locales y bloquea la rendición y el cobro hasta validar la evidencia.
  \item La exposición queda acotada a los datos ya presentes en el equipo y a un máximo de 14 horas.
\end{itemize}
